\newcommand{\zoomStartTime}{24.6}
\newcommand{\zoomEndTime}{25.3}
\newcommand{\resultFig}{figs/CONAC_figs}

%%%%%%%%%%%%%%%%%%%%%%%%%%%%%%%%%%%%%%%%%%%%%%%%%%%%%%%
% CHAPTER: HardNet++ on Neuro-Adaptive Control
%%%%%%%%%%%%%%%%%%%%%%%%%%%%%%%%%%%%%%%%%%%%%%%%%%%%%%%
\chapter{Hard-CONAC: Integration of HardNet++ and CONAC}
\label{CONAC:chap:hard:CONAC}

This chapter introduces the proposed Hard-Constrained Optimization-Based Neuro-Adaptive Control (Hard-CONAC) architecture, which integrates the HardNet++ framework with the CONAC approach.
Using the HardNet++ framework, the proposed Hard-CONAC architecture can handle hard constraints on the control input, which is a significant improvement over the original CONAC approach that only considers soft constraints.

The remainder of this chapter formulates the control problem in Section~\ref{CONAC:sec:problem:formulation}, introduces the DNN with a differentiable projection layer in Section~\ref{CONAC:sec:DNN:arch}, develops Hard-CONAC and its adaptation laws in Section~\ref{CONAC:sec:proposed:method}, analyzes its stability in Section~\ref{CONAC:sec:stability:analysis}, validates its performance in Section~\ref{CONAC:sec:validation}, and concludes in Section~\ref{CONAC:sec:conclusion}.

%%%%%%%%%%%%%%%%%%%%%%%%%%%%%%%%%%%%%%%%%%%%%%%%%%%%%%%
% PROBLEM FORMULATION
%%%%%%%%%%%%%%%%%%%%%%%%%%%%%%%%%%%%%%%%%%%%%%%%%%%%%%%
\section{Problem Formulation}
\label{CONAC:sec:problem:formulation}

Consider a nonlinear system with the following dynamics:
\begin{equation}
  \label{CONAC:eq:sys}
  \ddtt\mv{x}
  =
  \mv{f}(\mv{x})
  +
  \mm{B}(\mv{x})
  \mysat
  \left(
    \mv{u}
  \right)
  ,
\end{equation}
where $\mv{x} \in \R^{n}$ is the state vector, $\mv{u} \in \R^{m}$ is the control input, $\mv{f}:\R^n\to\R^m$ is an unknown system function, $\mm{B}:\R^n\to\R^{n\times m}$ is a unknown control effective matrix, and $\mysat(\cdot)$ represents the saturation function.

The system is subject to the following assumptions:

\begin{assum} 
  \label{CONAC:assum:B:bound}
  The control effectiveness matrix $\mm{B}(\mv{x})$ is nonsingular and bounded as   
  $
    \underline{b}\mm{I}_n
    \preceq 
    \mm{B}(\mv{x})
    \preceq 
    \bar{b}\mm{I}_n
    ,
  $
  where $\underline{b},\bar{b}\in\R_{>0}$ are unknown positive constants.
\end{assum}
\begin{assum}[{\cite[Assump. 1(b)]{Boulkroune:2010aa}}] 
  \label{CONAC:assum:B:bound:grad}
  There exists an unknown positive function $\rho_{\beta}(\mv{x}):\R^n\to\R_{>0}$ such that $\tfrac{1}{2}\mv{v}^\top\left(\ddtt\mm{B}(\mv{x})^{-1}\right)\mv{v} \le \rho_{\beta}(\mv{x})\norm{\mv{v}}^2$, $\forall \mv{v}\in\R^n$, and $\forall \mv{x}\in\R^n$. 
\end{assum}

The saturation function $\mysat(\cdot)$ may be chosen arbitrarily, provided that it maps the unconstrained input $\mv{u}$ into the admissible input set $\Omega_u\subset\R^n$.
The admissible input set is defined as follows:
\begin{equation}
  \label{CONAC:eq:adm:input:set}
  \Omega_u
  =
  \left\{
    \mv{u}\in\R^n
    \mid
    \underline{\mv{b}}
    \leq
    \mv{c}(\mv{u})
    \leq
    \overline{\mv{b}}
  \right\}
  ,
\end{equation}
where $c(\mv{u})$ is a constrained function vector and $\underline{b},\overline{b}\in\R^n$ are the lower and upper bounds of the constraints, respectively.

\par\medskip

Additionally, the desired trajectory $\mv{x}_d(t):\R_{\ge0}\to\R^n$ is assumed to be given and satisfies the following assumption.

\begin{assum} 
  \label{CONAC:assum:xd}
  The desired trajectory $\mv{x}_d:=\mv{x}_d(t)$ is sufficiently smooth, and $\mv{x}_d$ and $\ddtt\mv{x}_d$ are bounded.
\end{assum}

The control objective is to design a NAC for the unknown system \eqref{CONAC:eq:sys} such that $\mv{x}$ tracks the desired trajectory $\mv{x}_d$, the multiple convex input constraints defined by $\mysat(\cdot)$ are satisfied, and closed-loop stability is guaranteed.


%%%%%%%%%%%%%%%%%%%%%%%%%%%%%%%%%%%%%%%%%%%%%%%%%%%%%%%
% DNN ARCHITECTURE WITH DIFFERENTIABLE PROJECTION LAYER
%%%%%%%%%%%%%%%%%%%%%%%%%%%%%%%%%%%%%%%%%%%%%%%%%%%%%%%
\section{DNN Architecture with Differentiable Projection Layer (DPL)}
\label{CONAC:sec:DNN:arch}

\begin{figure}[t]
  \centering
  \includegraphics[width=0.99\linewidth]
  {
    figs/general_figs/ConstrainedDNN.drawio.pdf
  }
  \caption{
    Example illustration of the Hard-CONAC architecture whose node numbers of the hidden layers are all set to $3$ and the size of input and output are both $2$.
  }
  \label{CONAC:fig:hard:CONAC:arch}
\end{figure}

We first present the architecture of the DNN which is used in the proposed Hard-CONAC.
The architecture of the proposed Hard-CONAC is illustrated in Fig.~\ref{CONAC:fig:hard:CONAC:arch}.
As shown in the figure, the output of the unconstrained DNN is projected onto the feasible set defined by the constraints using the differentiable projection layer (DPL).

The unconstrained DNN $\NN_u(\NNin;\wth)$ is parameterized by the weight vector $\wth\in\R^{\Xi}$ and the input vector $\NNin\in\R^{l_0+1}$, where $l_0$ is the number of input nodes and $\Xi$ is the total number of weights in the DNN.
Specifically, the unconstrained DNN $\NN_u$ is recursively defined as follows:
\begin{equation}
  \label{CONAC:eq:DNN:uncon}
  \NN_i
  =
  \begin{cases}
    \wV_i^\top\act_i
    ,
    & 
    i \in \{1,\ldots,L\}
    \\
    \wV_0^\top\NNin
    ,
    &
    i = 0
    ,
  \end{cases}
\end{equation}
where $\wV_i\in\R^{l_{i}+1}\times l_{i+1}$ is the weight matrix of the $i$-th layer, and $\act_i\in\R^{l_{i}+1}$ is the activation vector of the $i$-th layer.
By adopting the differentiable projection layer (DPL) from HardNet++ proposed in \cite{Goertzen:2026aa}, the constrained DNN $\NN$ is defined as follows:
\begin{equation}
  \label{CONAC:eq:DNN:con}
  \NN
  =
  \Pi(\NN_u)
  ,
\end{equation}
where $\Pi:\R^{l_{L+1}}\to\R^{l_{L+1}}$ is the projection operator that projects the unconstrained DNN output $\NN_u$ onto the feasible set defined by the constraints.
The constraints are defined as follows:
\begin{equation}
  \label{CONAC:eq:cstr}
  \underline{\mv{b}}
  \leq
  \mv{c}(\NN_u)
  \leq
  \overline{\mv{b}}
  ,
\end{equation}
where $\mv{c},\underline{\mv{b}},\overline{\mv{b}}$ coincide with those in \eqref{CONAC:eq:adm:input:set} without their arguments for brevity.

The projection operator $\Pi(\NN_u)$ corrects $\NN_u$ to satisfy the constraints in \eqref{CONAC:eq:cstr} by minimizing the residual of the constraints, which is defined as follows:
\begin{equation}
  \label{CONAC:eq:proj:residual}
  \mv{r}(\NN_u) 
  = 
  \myReLU
  \left(
    \mv{c}(\NN_u) - \overline{\mv{b}}
  \right)
  +
  \myReLU
  \left(
    \underline{\mv{b}} - \mv{c}(\NN_u)
  \right)
  ,
\end{equation}
where $\mv{r}:=\mv{r}(\NN_u)$ is the residual of the constraints, and $\myReLU(\cdot):=\max(0, \cdot)$ is the rectified linear unit function.

\begin{algorithm}[!t]
  \caption{Differentiable Projection Layer (DPL)}
  \label{CONAC:alg:proj}
  \begin{algorithmic}[1]
    \STATE \textbf{Input:}
    Unconstrained DNN output $\NN_u$, lower and upper bounds $\mv{b}_{\ell}$ and $\mv{b}_{u}$, maximum number of iterations $K$, damping matrix $\mm{E}\succ\mm{0}$, and tolerance $\varepsilon_{\mathrm{tol}}$

    \STATE \textbf{Initialize:}
    $\mv{y}^{(0)}\leftarrow\NN_u$

    \FOR{$k=0$ to $K-1$}
      \STATE Evaluate the constraint function
      \[
        \mv{c}^{(k)}
        \leftarrow
        \mv{c}\bigl(\mv{y}^{(k)}\bigr)
      \]

      \STATE Compute the constraint residual
      \[
        \mv{r}^{(k)}
        \leftarrow
        \myReLU
        \left(\mv{c}^{(k)}-\overline{\mv{b}}\right)
        -
        \myReLU
        \left(\underline{\mv{b}}-\mv{c}^{(k)}\right)
      \]

      \IF{$\|\mv{r}^{(k)}\|_{\infty}
              \leq\varepsilon_{\mathrm{tol}}$}
        \STATE \textbf{break}
      \ENDIF

      \STATE Compute the constraint Jacobian
      \[
        \mm{J}^{(k)}
        \leftarrow
        \left.
        \frac{\partial\mv{c}(\mv{y})}
            {\partial\mv{y}}
        \right|_{\mv{y}=\mv{y}^{(k)}}
      \]

      \STATE Compute the damped correction
      \[
        \Delta\mv{y}^{(k)}
        \leftarrow
        -{\mm{J}^{(k)}}^{\top}
        \left(
          \mm{J}^{(k)}{\mm{J}^{(k)}}^{\top}
          +\mm{E}
        \right)^{-1}
        \mv{r}^{(k)}
      \]

      \STATE Update the projected variable
      \[
        \mv{y}^{(k+1)}
        \leftarrow
        \mv{y}^{(k)}+\Delta\mv{y}^{(k)}
      \]
    \ENDFOR

    \STATE \textbf{Output:}
    \[
      \Pi(\NN_u)
      \leftarrow
      \mv{y}^{(K)}
    \]
  \end{algorithmic}
\end{algorithm}

\par\medskip

The projection operator $\Pi(\NN_u)$ iteratively linearize the constraint function $\mv{c}$ and projects the input onto the feasible set defined by the linearized constraints.
For better intuition, the algorithm of the projection operator $\Pi(\cdot)$ is summarized in Algorithm~\ref{CONAC:alg:proj}.

When the unconstrained DNN output $\NN_c$ violates any constraint, the DPL initializes $\mv{y}^{(0)}=\NN_c$ and constructs a linearized projection step by locally approximating the constraint functions.
For instance, at step $k$, the linearized constraints are constructed at $\mv{y}^{(k)}$, as follows:
\begin{equation}
  \underline{\mv{b}}
  \le
  c(\mv{y}^{(k)}) + \mm{J}^{(k)}(\mv{y}^{(k+1)}-\mv{y}^{(k)})
  \le
  \overline{\mv{b}}
  ,
\end{equation}
where $\mm{J}^{(k)}$ is the Jacobian of the constraint function evaluated at $\mv{y}^{(k)}$.
Invoking the linearized constraints, HardNet-Aff presented in \cite{Min:2025aa} 
\cred{
  Need to be added!
}

Then, by adding a damping term $\mm{E}=\mm{E}^\top\succ\mm{0}$, the projection step is computed as follows:
\begin{equation}
  \label{CONAC:eq:proj:step}
  \mv{y}^{(k+1)}
  =
  \mv{y}^{(k)}
  -
  {\mm{J}^{(k)}}^{\top}
  \left(
    \mm{J}^{(k)}{\mm{J}^{(k)}}^{\top}
    +\mm{E}
  \right)^{-1}
  \mv{r}^{(k)}
  .
\end{equation}
Finally, when the infinite norm of the constraint residual $\mv{r}^{(k)}$ is less than a specified tolerance $\varepsilon_{\mathrm{tol}\in\R_{>0}}$, the DPL terminates the iteration and outputs the projected variable $\mv{y}^{(K)}$ as the final output of the projection operator $\Pi(\NN_u)$.


%%%%%%%%%%%%%%%%%%%%%%%%%%%%%%%%%%%%%%%%%%%%%%%%%%%%%%%
% PROPOSED METHOD
%%%%%%%%%%%%%%%%%%%%%%%%%%%%%%%%%%%%%%%%%%%%%%%%%%%%%%%
\section{Proposed Method}
\label{CONAC:sec:proposed:method}

\begin{figure}[t]
  \centering
  % \fbox{
  %   \begin{minipage}[c][0.2\textheight][c]{0.95\linewidth}
  %     \centering
  %     \textbf{Will be updated}
  %   \end{minipage}
  % } 
  \includegraphics[width=0.8\linewidth]
  {
    figs/general_figs/HardCONAC.drawio.pdf
  }
  \caption{
  Overall architecture of the proposed Hard-CONAC.
  }
  \label{CONAC:fig:proposed:arch}
\end{figure}


In this section, we present the proposed Hard-CONAC architecture, which integrates the HardNet++ framework with the CONAC approach.
The overall architecture of the proposed Hard-CONAC is illustrated in Fig.~\ref{CONAC:fig:proposed:arch}.

% —————————————————————————————————————————————————————
% PROPOSED METHOD: NAC formulation
% —————————————————————————————————————————————————————
\subsection{Neuro-Adaptive Control Formulation}

In this section, the NAC formulation is presented.
First, the error dynamics $\mv{e}:=\mv{x}-\mv{x}_d$ is defined, and the error dynamics is rewritten as follows:
\begin{equation}
  \label{CONAC:eq:err:dyn}
  \ddtt\mv{e}
  =
  \mv{f}(\mv{x}) + \mm{B}(\mv{x})\mysat(\mv{u}) - \ddtt\mv{x}_d
  .
\end{equation}

Then, consider the following Lyapunov function candidate:
\begin{equation}
  \label{CONAC:eq:Vc}
  V_c = \tfrac{1}{2}\mv{e}^\top\mm{B}^{-1}(\mv{x})\mv{e}
  .
\end{equation}
Invoking Assumptions~\ref{CONAC:assum:B:bound:grad}, the time derivative of $V_c$ can be bounded as follows:
\begin{align} \label{CONAC:eq:Vc:dot:1} \nonumber
  \ddtt V_c
  =&
  \tfrac{1}{2}\mv{e}^\top
  \left(
    \ddtt\mm{B}^{-1}
  \right)
  \mv{e}
  +
  \mv{e}^\top\mm{B}^{-1}
  \left(
    \mv{f} + \mm{B}\mysat(\mv{u}) 
    - \ddtt\mv{x}_d
  \right)
  \\
  \le&
  \rho_{\beta}(\mv{x})
  \norm{\mv{e}}^2
  +
  \mv{e}^\top\mv{u}
  +
  \mv{e}^\top\mm{B}^{-1}
  \left(
    \mv{f}- \ddtt\mv{x}_d
  \right)
  +
  \mv{e}^\top
  \Delta\mv{u}
  ,
\end{align}
where the arguments of $\mv{f}(\mv{x})$ and $\mm{B}(\mv{x})$ are omitted for brevity and $\Delta\mv{u} := \mysat(\mv{u}) - \mv{u}$.

Notably, as shown in the previous section, the proposed Hard-CONAC architecture ensures that the control input $\mv{u}$ always satisfies the input constraints, resulting in 
\begin{equation}
  \Delta\mv{u} = \mv{0}
  .
\end{equation}
Therefore, the time derivative of the Lyapunov function candidate $V_c$ can be further simplified as follows:
\begin{equation} \label{CONAC:eq:Vc:dot:2}
  \ddtt V_c
  \le
  \rho_{\beta}(\mv{x}) + \mv{e}^\top \mv{u}
  +
  \mv{e}^\top\mm{B}^{-1}
  \left(
    \mv{f}- \ddtt\mv{x}_d
  \right)
  .
\end{equation}

A desired control law $\mv{u}_d$ such that the tracking error $\mv{e}$ is stable in the absence of input saturation is selected as 
\begin{equation} \label{eq:ctrl:ideal}
  \mv{u}_d
  :=
  -
  \mm{P}
  \mv{e}
  - 
  \rho_{\beta}(\mv{x})
  \mv{e}
  +
  \mm{B}^{-1}
  \left(
    - \mv{f} 
    + \ddtt\mv{x}_d
  \right)
  % +
  % h(\mv{x},\mv{x}_d)
  ,
\end{equation}
where $\underline{p}\mm{I}_n \preceq \mm{P}=\mm{P}^\top \preceq \overline{p}\mm{I}_n$ for known positive constants $\underline{p},\overline{p}\in\R_{>0}$.
% The function $h:\R^n\times\R^n\to\R^n$ is correction term to satisfy the ideal control law $\mv{u}^*$ within the input constraints, \ie $\mv{u}^*\in\Omega_u$.

By substituting $\mv{u}_d$ in \eqref{eq:ctrl:ideal} into \eqref{CONAC:eq:Vc:dot:3}, we have
\begin{align} \label{CONAC:eq:Vc:dot:3} 
  \ddtt V_c
  \le &
  -\mv{e}^\top
  \mm{P}
  \mv{e}
  +
  \mv{e}^\top
  \left(
    \mv{u} - \mv{u}_d
  \right)
  \\ \nonumber
  \le & 
  -
  \underline{p}\norm{\mv{e}}^2
  +
  \mv{e}^\top
  \left(
    \mv{u} - \mv{u}_d       
  \right)
  .
\end{align}
Notably, in the ideal case, \ie $\mv{u} = \mv{u}_d$, the tracking error $\mv{e}$ asymptotically converges to zero with decay rate $\underline{p}\overline{b}$, since 
$
  \tfrac{1}{\overline{b}}
  \ddtt
  \left(
    \norm{\mv{e}}^2
  \right)
  \le
  \ddtt V_c 
  \le 
  -\underline{p}\norm{\mv{e}}^2
$.

Since the ideal control input $\mv{u}_d$ is unknown due to the unknown system dynamics, we approximate $\mv{u}_d$ using the DNN architecture with the DPL which is defined in the previous section, as follows:
\begin{equation}
  \label{CONAC:eq:ctrl:ideal}
  \mv{u}_d
  =
  \NN(\NNin;\idealwth)
  +
  \mv{\epsilon}
  ,
\end{equation}
where $\idealwth\in\R^{\Xi}$ is the ideal weight vector, and $\mv{\epsilon}\in\R^n$ is the approximation error of the DNN.
The ideal weight vector and the approximation error are bounded for a compact set $\Omega_{\NNin}\subset\R^{l_0+1}$ which is also known as the approximation region, as follows:
\begin{equation}
  \norm{\idealwth} 
  \le 
  \overline{\theta}
  , 
  \qquad 
  \norm{\mv{\epsilon}} 
  \le 
  \overline{\epsilon}
  , 
  \qquad 
  \forall 
  \NNin\in\Omega_{\NNin}
  ,
\end{equation}
where $\overline{\theta},\overline{\epsilon}\in\R_{>0}$ are unknown positive constants.

As the ideal weight vector $\idealwth$ is unknown, we estimate it using the weight vector $\estwth\in\R^{\Xi}$, which is updated online using the adaptation law derived in the subsequent section.
The actual control input $\mv{u}$ is then defined as follows:
\begin{equation}
  \label{CONAC:eq:ctrl:actual}
  \mv{u}
  =
  \NN(\NNin;\estwth)
  =
  \Pi\left(\NN_u(\NNin;\estwth)\right)
  .
\end{equation}

By substituting \eqref{CONAC:eq:ctrl:ideal} and \eqref{CONAC:eq:ctrl:actual} into \eqref{CONAC:eq:Vc:dot:3}, we have
\begin{align} \label{CONAC:eq:Vc:dot:4}
  \ddtt V_c
  \le &
  -\underline{p}\norm{\mv{e}}^2
  +
  \mv{e}^\top
  \left(
    \NN(\NNin;\estwth)
    -
    \NN(\NNin;\idealwth)
    -
    \mv{\epsilon}
  \right)
  \\ \nonumber
  \le &
  -\underline{p}\norm{\mv{e}}^2
  +
  \mv{e}^\top
  \pptfrac{\estNN}{\estwth}\tilde{\wth}
  +
  \mv{e}^\top
  \underbrace{
  \mathcal{O}\left(\norm{\tilde\wth}\right)
  +
  \mv{\epsilon}
  }_{=:\mv{\Delta}}
  ,
\end{align}
where $\tilde{\wth}:=\estwth-\idealwth$ is the weight estimation error, and $\pptfrac{\estNN}{\estwth}$ is the Jacobian of the constrained DNN with respect to the weight vector $\estwth$.
The high-order term is denoted by $\mathcal{O}(\norm{\tilde\wth})$ and is bounded only if $\norm{\tilde\wth}$ is bounded.
The term $\mv{\Delta}$ includes the high-order term and the approximation error.

% —————————————————————————————————————————————————————
% PROPOSED METHOD: Adaptation law derivation
% —————————————————————————————————————————————————————
\subsection{Adaptation Law Derivation}

The constrained control problem stated in Section~\ref{CONAC:sec:problem:formulation} can be formulated as a constrained optimization problem, which is defined as follows:
\begin{subequations} 
  \label{CONAC:eq:opt:prob} 
  \begin{align} 
    \min_{\estwth} \quad 
    & 
    J(\mv{e};\estwth)
    = 
    \tfrac{1}{2} \mv{e}^\top \mv{e} 
    +
    \tfrac{1}{2} \mv{r}^\top \mv{r}
    \label{CONAC:eq:opt:prob:obj} 
    \\ 
    \text{subject to} \quad  
    & 
    c^{\theta}_i(\estwth_i)=\tfrac{1}{2}
    \left(
      \norm{\estwth_i}^2
      -
      \overline{\theta}_i^2 
    \right)
    \le 0, 
    \qquad
    \forall i\in\{0,1,\ldots,L\}
    ,
    \label{CONAC:eq:opt:prob:const1} 
    \\ 
    & 
    \underline{\mv{b}}(\NNin)
    \leq
    \mv{c}^u(\NN)
    \leq
    \overline{\mv{b}}(\NNin)
    ,
    \label{CONAC:eq:opt:prob:const2} 
  \end{align}%
\end{subequations}
where $J(\mv{e};\estwth)$ is the objective function, and $c^{\theta}_i(\estwth_i)$ is the constraint function for the weight vector $\estwth_i$ of the $i$-th layer.

The objective function includes the instantaneous tracking error $\mv{e}$ and the constraint residual $\mv{r}$, which is defined in \eqref{CONAC:eq:proj:residual}.

\cred{...}

The adaptation law for the adaptive variables, $\estwth$, and $[\lambda_i]_{j\in\{0,1,\ldots,L\}}$, is derived to solve \eqref{eq:opt:prob:saddle} by using a primal-dual gradient update law, as follows:
\begin{subequations} 
  \label{CONAC:eq:adap}
  \begin{align}
    \ddtt{\estwth}
    &
    = 
    -\alpha \pptfrac{L}{\estwth}
    =
    -\alpha 
    \left(
      \pptfrac{J}{\estwth}
      +
      \textstyle\sum_{i=0}^L
      \lambda_i 
      \estwth_i
    \right),
    \label{CONAC:eq:adap:th}
    \\
    \ddtt\lambda_i
    & 
    = 
    \myproj_{\lambda_i\ge0}
    \left(
      \beta_i
      \pptfrac{L}{\lambda_i}
    \right)
    =
    \myproj_{\lambda_i\ge0}
    \left(
      \beta_i
      c^{\theta}_{i} 
    \right)
    ,
    \qquad \forall i\in\{0,1,\ldots,L\},
    \label{CONAC:eq:adap:lbd}
  \end{align}%
\end{subequations}

% —————————————————————————————————————————————————————
% PROPOSED METHOD: FORWARD-MODE AUTOMATIC DIFFERENTIATION
% —————————————————————————————————————————————————————
\subsubsection{Forward-Mode Automatic Differentiation for DPL}

As shown in the previous section, it is not straightforward to compute the gradient of the constrained DNN $\NN$ with respect to its input, \ie the unconstrained DNN $\NN_u$, due to the projection operator $\Pi(\cdot)$.
We adopt the forward-mode automatic differentiation technique \cite[Chap. 8.2]{Nocedal:2006aa} to compute the gradient of the each projection layers in the DPL.

First, 

%%%%%%%%%%%%%%%%%%%%%%%%%%%%%%%%%%%%%%%%%%%%%%%%%%%%%%%
% STABILITY ANALYSIS
%%%%%%%%%%%%%%%%%%%%%%%%%%%%%%%%%%%%%%%%%%%%%%%%%%%%%%%
\section{Stability Analysis}
\label{CONAC:sec:stability:analysis}

In this section, the closed-loop stability under the proposed Hard-CONAC is analyzed.

\begin{theorem} \label{thm:stability}
  \cred{add theorem statement}
  % Consider the closed-loop system consisting of \eqref{eq:sys}, the control law \eqref{eq:ctrl:NN:actual}, and the adaptation law \eqref{eq:adap}.
  % Suppose that Assumptions~\ref{assum:B:bound}--\ref{assum:xd} hold and the initial condition belongs to the compact invariant set 
  % $
  % \prod_{\ell=0}^k \mathcal{S}_\ell(\overline V_\ell)
  % $ 
  % constructed in \eqref{eq:Sk:def} and \eqref{eq:Si:def}.
  % Then, the tracking error $\mv e$, the DNN weights $\estwth$, and the Lagrange multipliers $\sigma_i,\ \forall i \in \mathcal{K},$ and $\lambda_j,\ \forall j \in \mathcal{I},$ are uniformly ultimately bounded.
\end{theorem}

\begin{proof}

Consider the following Lyapunov function candidate:
\begin{equation}
  \label{CONAC:eq:V}
  V
  =
  V_c
  +
  \tfrac{1}{2\alpha}\tilde{\wth}^\top\tilde{\wth}
  +
  \textstyle\sum_{i=0}^L
  \tfrac{1}{\beta_i}
  \left(
    \lambda_i - \lambda_i^o
  \right)^2
  .
\end{equation}
Taking the time derivative of $V$ yields
\begin{align} \label{CONAC:eq:V:dot:1}
  \ddtt V
  = &
  -\underline{p}\norm{\mv{e}}^2
  +
  \mv{e}^\top
  \left(
    \pptfrac{\estNN}{\estwth}\tilde{\wth}
    +
    \mv{\Delta}
  \right)
  -
  \tilde{\wth}^\top
  \left(
    \pptfrac{\estNN}{\estwth}^\top
    \mv{e}
    +
    \left(
      \mm{D}^{(0)}
      \mm{J}^{(0)}
      \pptfrac{\estNN_u}{\estwth}
    \right)^\top
    \mv{r}
    +
    \mydiag(\mv{\lambda})\estwth
  \right)
  +
  \sum_{i=0}^L
  \left(
    \lambda_i - \lambda_i^o
  \right)
  \left(
    \estwth_i^\top
    \estwth_i
    -
    \overline{\theta}_i^2
  \right)
  \\
  =
  &
  -\underline{p}\norm{\mv{e}}^2
  +
  \mv{e}^\top
  \mv{\Delta}
  -
  \tilde{\wth}^\top
  \left(
    \mm{D}^{(0)}
    \mm{J}^{(0)}
    \pptfrac{\estNN_u}{\estwth}
  \right)^\top
  \mv{r}
  +
  \sum_{i=0}^L
  \left(
    -
    \lambda_i
    \tilde{\wth}^\top
    \estwth
    +
    \left(
      \lambda_i - \lambda_i^o 
    \right)
    \left(
      \estwth_i^\top
      \estwth_i
      -
      \overline{\theta}_i^2
    \right)
  \right)
\end{align}


\end{proof}

%%%%%%%%%%%%%%%%%%%%%%%%%%%%%%%%%%%%%%%%%%%%%%%%%%%%%%%
% VALIDATION
%%%%%%%%%%%%%%%%%%%%%%%%%%%%%%%%%%%%%%%%%%%%%%%%%%%%%%%
\section{Validation}
\label{CONAC:sec:validation}



\begin{table}[t]
  \renewcommand{\arraystretch}{1.3}
  \caption{Controllers and parameters used for comparative analysis.}
  \label{table:controller}
  \centering
  % \begin{tabular}{m{.6cm} m{4cm} m{2.7cm}}
  \begin{tabular}{m{.6cm} m{6cm} m{4.7cm}}
  % \begin{tabular}{c l l}
  \hline
    &\bf{Description and parameters}&\bf{Input constraint}\\
  \hline
  \hline
    \multirow{3}{*}{
      \shortstack{
        {(C$_1$)}
        \\
        \protect\colorLine{blue}{solid}
      }
    }& Hard-CONAC 
    & 
      \multirow{3}{*}{
        \shortstack[l]{
          $-\overline{\tau}_2 \le \tau_2 \le \overline{\tau}_2$, 
          \\
            $\norm{\mv{\tau}} \le \overline{\tau}$,
          \\
            ($\overline{\tau}_2=3.8, \overline{\tau}=11$)
        }
      }
    \\
    & $\epsilon=0 $ & 
    \\
    & 
    $\overline{\theta}_i=15$
    &
    \\
  \hline
    \multirow{3}{*}{
      \shortstack{
        {(C$_2$)}
        \\
        \protect\colorLine{my_cyan}{solid}
      }
    }& CONAC &  
    \\
    & $\beta^{u}_{\overline{\tau}_2}=\beta^{u}_{\underline{\tau}_2}=1000$, $\beta^{u}_{\overline{\tau}}=100,$ & same as (C$_1$) 
    \\
    & 
    $\overline{\theta}_i=15,\beta_i^\theta=1,\forall i\in\{0,1,2\}$,
    &
    \\
  \hline
    \multirow{2}{*}{
    \shortstack{  
      (C$_3$) 
      \\
      \protect\colorLine{gray}{solid}
    }
    }& NAC without input constraints& 
    \multirow{2}{*}{not considered}
    \\
    & $\overline{\theta}=15$ \cite{Patil:2022aa} & 
    \\  
\hline
\hline
    \multicolumn{3}{m{12.1cm}}{
      DNN and other parameters for all controllers:
      $
        \alpha=0.2
        ,
        \mm{P}=\mm{I}_2
        ,
      $
    }
    \\
    \multicolumn{3}{m{12.1cm}}{
      $
        \NNin 
        \!=\!
        \left(
          (\mv{q}-\mv{q}_d)^\top,
          \ddtt(\mv{q}-\mv{q}_d)^\top,
          \tfrac{\der^2}{\der t^2}\mv{q}_d^\top,
          1
        \right)^\top
        ,
        \mm{\Lambda} = \mydiag(6,13)
        ,
      $
    }
    \\
    \multicolumn{3}{m{12.1cm}}{
      $
        l_0=6,l_1=4,l_2=4,l_3=2
        ,
        \estwth(0)\sim\mathcal{U}\bigl([-0.5,0.5]^{\Xi}\bigr)
      $
    }
    \\
\hline
  \end{tabular}
\end{table}

\begin{figure*}[t]
  \centering
  \subfloat[Tracking performance of $1$st joint angle $q_1$.]{
    \includegraphics[width=.49\linewidth]
    {
      \resultFig/Fig1.eps
    }%
    \label{fig:ctrl:result:general:q1}}
  \hfill
  \subfloat[Tracking performance of $2$nd joint angle $q_2$.]{
    \includegraphics[width=.49\linewidth]
    {
      \resultFig/Fig2.eps
    }%
    \label{fig:ctrl:result:general:q2}}
  \vfill
  % \vfill
  \subfloat[Commanded control input of the $1$st joint $\tau_1$.]{
  \includegraphics[width=.49\linewidth]
    {
      \resultFig/Fig5.eps
    }%
    \label{fig:ctrl:result:general:tau:1}}
  \hfill
    \subfloat[Commanded control input of the 2nd joint $\tau_2$.]{
    \includegraphics[width=.49\linewidth]
    {
      \resultFig/Fig6.eps
    }%
    \label{fig:ctrl:result:general:tau:2}}
  \vfill
  \subfloat[Norm of the commanded control input $\mv{\tau}$.]{
    \includegraphics[width=.49\linewidth]
    {
      \resultFig/Fig7.eps
    }%
    \label{fig:ctrl:result:general:tau:norm}}
  \hfill
  % \vfill
  \subfloat[Norm of DNN weights ${\estwth}_i$.]{
    % \includegraphics
    \includegraphics[width=.49\linewidth]
    {
      \resultFig/Fig10.eps
    }%
    \label{fig:ctrl:result:general:weight}}        
  \caption{
      Real-time implementation results of (C$_1$) \protect\colorLine{blue}{solid}, (C$_2$) \protect\colorLine{my_cyan}{solid}, (C$_3$) \protect\colorLine{gray}{solid}, and desired trajectory $\qd$ \protect\colorLine{red}{dashed}.
      % \vspace{-5mm}
  }
\label{fig:ctrl:result:general}
\end{figure*}






\begin{figure*}[t]
  \centering
  \subfloat[Number of the linearized projection steps in the DPL.]{
    \includegraphics[width=0.48\linewidth]
    {
      \resultFig/Fig11.eps
    }
    \label{fig:ctrl:result:detail:q:top:view}
  }
  \hfill
  \subfloat[Convergence of the DPL.]{
    \includegraphics[width=0.49\linewidth]
    {
      \resultFig/Fig13.eps
    }
    \label{fig:ctrl:result:detail:u:top:view}
  }
  \caption{  
    Detail results of the control implementation of (C$_1$).
  }
  \label{fig:ctrl:result:detail}
\end{figure*}


\begin{figure*}[t]
  \centering
  \subfloat[Trajectory of the tracking error of each joint angle $\mv{q}$ during the time interval {$[\zoomStartTime, \qty{\zoomEndTime}{\second}]$}.]{
    \includegraphics[width=0.48\linewidth]
    {
      \resultFig/Fig14.eps
    }
    \label{fig:ctrl:result:detail:q:top:view}
  }
  \hfill
  \subfloat[Trajectory of the commanded control input $\mv{\tau}$ during the time interval {$[\zoomStartTime, \qty{\zoomEndTime}{\second}]$}.]{
    \includegraphics[width=0.48\linewidth]
    {
      \resultFig/Fig8.eps
    }
    \label{fig:ctrl:result:detail:u:top:view}
  }
  \vfill
  \subfloat[Number of the linearized projection steps in the DPL of (C$_1$) during the time interval {$[\zoomStartTime, \qty{\zoomEndTime}{\second}]$}.]{
    \includegraphics[width=0.48\linewidth]
    {
      \resultFig/Fig12.eps
    }
    \label{fig:ctrl:result:detail:multipliers}
  }
  \hfill
  \subfloat[Lagrange multiplier associated with the power source limit constraint for (C$_1$) during the time interval {$[\zoomStartTime, \qty{\zoomEndTime}{\second}]$}.]{
    \includegraphics[width=0.48\linewidth]
    {
      \resultFig/Fig9.eps
    }
    \label{fig:ctrl:result:detail:multipliers}
  }
  \caption{  
    Zoomed-in view of the tracking error, commanded control input, and
Lagrange multiplier $\lambda_{\overline{\tau}}$ during the time interval
$[\zoomStartTime, \qty{\zoomEndTime}{\second}]$ of (C$_1$) \protect\colorLine{blue}{solid}, (C$_2$) \protect\colorLine{my_cyan}{solid}, and (C$_3$) \protect\colorLine{gray}{solid}.
  }
  \label{fig:ctrl:result:detail}
\end{figure*}

%%%%%%%%%%%%%%%%%%%%%%%%%%%%%%%%%%%%%%%%%%%%%%%%%%%%%%%
% CONCLUSION
%%%%%%%%%%%%%%%%%%%%%%%%%%%%%%%%%%%%%%%%%%%%%%%%%%%%%%%
\section{Conclusion}
\label{CONAC:sec:conclusion}


